\section{Introduction}

Schema drift is still found too late in many data systems. A producer renames a field, reorders a record, widens optionality, or changes the nullability of a nested collection, and the break often appears only when a job touches real data. Spark offers runtime schema comparators and schema-aware readers, but those checks still run after the program has been built and after an external boundary has already been crossed. A pipeline can therefore look typed in code while the key compatibility question is deferred to execution time~\cite{spark_datatype}.

Typed-Dataset layers and table-level schema enforcement each close part of this gap but leave a seam. Typed-Dataset libraries such as \texttt{frameless} and \texttt{iskra} require that raw Spark \texttt{DataFrame} code be replaced with a typed wrapper across the job, which raises adoption cost in practice~\cite{frameless,iskra}. Table-level schema enforcement such as Delta Lake and Apache Iceberg operates at write time against a stored schema rather than at compile time against a declared Scala type, and therefore still runs after job submission~\cite{delta_lake,apache_iceberg}. Neither layer currently provides compile-time proof that a declared producer type conforms to a declared contract type under an explicit policy before the job is submitted at all. That is the seam this paper's artifact targets. The value of closing it is argued here and left for future work to measure.

For typed JVM pipelines, much of that question is structural and available earlier. If the producer and sink contract are both represented as Scala types, then field names, nesting, collection shape, optionality, and defaulted fields are visible to the compiler. Earlier Scala derivation routes such as \texttt{shapeless} and Scala~3 \texttt{Mirror}-based derivation show adjacent implementation paths, but this artifact uses quoted reflection directly because it exposes \texttt{TypeRepr}, symbols, and type structure without introducing a separate generic wrapper layer~\cite{shapeless,scala3_reflection}. That gives us a concrete way to compute a normalized structural model at compile time and reject incompatible producer-to-contract pairs before a pipeline ships.

Compile time is not enough. Even if application code proves that a declared producer type conforms to a contract, the actual runtime schema can still drift at the data boundary. CSV, JSON, and other external inputs do not obey the Scala types used in pipeline code. The artifact in this paper therefore combines two checks: compile-time proof that the declared producer type conforms to the target contract under an explicit policy, and a runtime pin that validates the actual Spark \texttt{DataFrame} schema before write. The first shifts structural drift detection into compilation and CI; the second protects the external boundary where runtime data can still violate assumptions.

This paper presents a deliberately narrow framework for that combination. It derives normalized structural shapes from Scala~3 case classes, materializes compile-time evidence \texttt{SchemaConforms[Out, Contract, P]} under an explicit policy family, derives Spark schemas from the same contract types, and enforces a policy-aware runtime comparator at the sink boundary. Where Spark's documented comparators stop, the runtime layer adds a deeper check for nested collection optionality, which the artifact shows as a real drift source in tests~\cite{spark_datatype}.

The artifact is structural rather than semantic. It does not model business rules, temporal properties, cross-record invariants, or external schema-registry workflows. That boundary matters. Prior work on software contracts shows that contracts are useful in practice, but also that the space is much richer than boundary shape checking; temporal and effectful contracts, for example, are well beyond the scope of this artifact~\cite{contracts_in_practice,trace_contracts,effectful_software_contracts}.

Within that scope, this paper makes three contributions:

\begin{enumerate}[leftmargin=1.5em]
  \item A Scala~3 macro derivation of normalized structural shapes that rejects incompatible producer-to-contract pairs at compile time under an explicit, tested policy family, with path-rich drift diagnostics.
  \item A sink-boundary design that fuses that compile-time witness with a policy-aware Spark runtime comparator, including subset semantics for \texttt{Backward}/\texttt{Forward} and a nested-collection-optionality check that Spark's built-in comparators omit.
  \item A reproducible artifact demonstrating that this combination can be carried at low additional compile cost and at write-time runtime cost that is higher than Spark's built-in comparators but still within a single-digit microsecond per schema comparison.
\end{enumerate}

The paper is therefore best read as a mechanism paper with an honest artifact, not as a production-effectiveness paper. What is proven here is that a focused class of structural drift checks can move from runtime failure toward compile-time proof and CI while still retaining a runtime guard at the actual data boundary.
